\documentclass{article}
\usepackage{amsmath, amssymb, amsthm, graphicx}
\usepackage{pgfplots, tikz, cancel, framed, hyperref}

\graphicspath{{./Images}}


% \geometry{margin=1in}
\usepackage[left=1in, right=1in, top=0.75in, bottom=1in]{geometry}
\pgfplotsset{compat=1.18}
\setlength{\parskip}{0.8em}
\setlength{\parindent}{0pt}

\newtheoremstyle{solution}
  {15pt}
  {15pt}
  {}
  {}
  {\bfseries}
  {: }
  { }
  {}

\theoremstyle{solution}
\newtheorem*{solution}{SOLUTION}

\usetikzlibrary{calc,shapes}

\newcommand{\tikzmark}[1]{\tikz[overlay,remember picture] \node (#1) {};}
\pgfmathdeclarefunction{erf}{1}{%
  \pgfmathparse{(2/sqrt(pi))*((#1) - ((#1)^3)/3 + ((#1)^5)/10 - ((#1)^7)/42 + ((#1)^9)/216)}%
}


\pgfmathdeclarefunction{erfi}{1}{%
  \pgfmathparse{(2/sqrt(pi))*((#1) + ((#1)^3)/3 + ((#1)^5)/10 + ((#1)^7)/42 + ((#1)^9)/216)}%
}



\pgfplotsset{my style/.append style={axis x line=middle, axis y line=
middle, xlabel={$x$}, ylabel={$y$} }}

%------------------------
\title{TEMPLATE TITLE}
\author{Brendan Tea}
\date{\today}

\begin{document}
\maketitle

\section{Jamiiru Homework for Chapter 1}

\begin{itemize}
    \item 
\end{itemize}


\begin{enumerate}
    \item 
\end{enumerate}


wow this is so cool!! $$ $$, $ax^2+bx+c=0$, where


\section{Vector Calculus is all about this}
Generalized Stokes Theorem
\begin{equation}
    \int_{\omega}^{} d \omega = \int_{\partial \Omega}^{} \omega
\end{equation}

You are integrating over the boundary over the region.

\end{document}